\begin{definition}[3.2.4]
\label{IV.2.3.2.4-fr}
Sean $X$ un preesquema localmente noetheriano, $\mathcal F$ un
$\mathcal O_X$-módulo cuasicoherente. Se dice que $\mathcal F$ es
irredundante si $\operatorname{Ass}(\mathcal F)$ se reduce a un único
elemento $x$~; si $\mathcal F$ es de tipo finito, se dice que $\mathcal F$ es
íntegro si además $\mathcal F$ es reducido \textup{(es decir, si
$\operatorname{long}(\mathcal F_x)=1$)}. Se dice que un sub-$\mathcal O_X$-módulo
cuasicoherente $\mathcal G$ de un $\mathcal O_X$-módulo cuasicoherente
$\mathcal F$ es primario en $\mathcal F$ si
$\mathcal F/\mathcal G$ es irredundante.
\end{definition}

Para un esquema afín $X$, esta definición da, en particular, la noción de
\emph{módulo íntegro} sobre un anillo noetheriano $A$~; se dice que un
$A$-módulo $M$ es \emph{íntegro} si es de tipo finito, si $M$ es
\emph{primario} (es decir, si $\operatorname{Ass}(M)$ se reduce a un único
ideal primo $\mathfrak p$) y si además
$\operatorname{long}_{A_\mathfrak p}(M_\mathfrak p)=1$. Volviendo al caso de un
preesquema localmente noetheriano cualquiera $X$ y de un
$\mathcal O_X$-módulo coherente $\mathcal F$, se dice que $\mathcal F$ es
\emph{íntegro en un punto} $x\in X$ si $\mathcal F_x$ es un
$\mathcal O_x$-\emph{módulo íntegro}~: esto significa también que $x$
\emph{pertenece a un único ciclo primo asociado (necesariamente no
inmerso) a $\mathcal F$ y que $\operatorname{long}(\mathcal F_z)=1$ en su
punto genérico $z$}. Es claro que, si $\mathcal F$ es un
$\mathcal O_X$-módulo coherente íntegro, es íntegro en todo punto de $X$~;
recíprocamente, si $\mathcal F$ es \emph{íntegro en un punto} $x$, existe una
vecindad abierta $U$ de $x$ tal que $\mathcal F|U$

\oldpage[IV-2]{42}
sea un $\mathcal O_U$-módulo \emph{íntegro}~: basta, en efecto, tomar
$U$ de modo que $\mathcal F|U$ sea un $\mathcal O_U$-módulo reducido (3.2.2).

Se dice que el preesquema localmente noetheriano $X$ es \emph{irredundante} si
$\mathcal O_X$ es irredundante (lo que implica que $X$ es
\emph{irreducible})~; para que $X$ sea \emph{íntegro}, es necesario y suficiente
que el $\mathcal O_X$-módulo $\mathcal O_X$ sea íntegro ((\textbf I, 2.1.8) y
(3.2.1)). Si $\mathcal O_X$ es íntegro en un punto $x$, es decir, si
el anillo $\mathcal O_x$ es íntegro, se dice que $X$ es \emph{íntegro en el
punto $x$}. Se dice que un subpreesquema cerrado $Y$ de $X$ es \emph{primario}
en $X$ si el ideal $\mathcal J$ de $\mathcal O_X$ que define $Y$ es primario
en $\mathcal O_X$.
